% P1 section 5.6: editorial completion of the frozen closure/quotient clauses.
\section{Equivalence closure and quotient constructions}
\label{sec:p01-equivalence-closure-and-quotients}

To require the least equivalence relation containing a relation is not yet
to construct one. The following construction supplies existence and
minimality, and then shows exactly when a function can be defined on the
resulting classes. It completes the equivalence-closure definition and
quotient method of \cref{sec:strings-05,sec:set-prerequisites-04}
\cite[sections ``Relations and closures'' and ``Equivalence relations and
quotients'']{FrozenFormalLogicFoundationsClosure}.
The proofs and finite example here are editorial completions, not proofs
attributed to an unread external presentation.

All carriers are sets. A binary relation on \(X\) is a subset of
\(X\times X\); comparison of relations below is inclusion on that same
carrier. Neither finiteness of \(X\) nor a well-order of \(X\) is assumed.
The natural numbers \(\N\) include zero. Relations such as \(R\) and
\(E\) are mathematical relations between elements, not object-language
formulas or a new inference calculus.

\subsection{Finite zig-zags construct the closure}

\begin{hbdefinition}{def:p01-ec-zigzag}{Converse, steps and paths}
For \(R\subseteq X\times X\), its converse is
\[
 R^{-1}\defeq\{\langle y,x\rangle\in X\times X:xRy\}.
\]
This is relation converse, not the assertion that a function has an inverse.
A \emph{zig-zag path of length \(n\)} from \(x\) to \(y\) is a
finite sequence \(x_0,\ldots,x_n\) of elements of \(X\), where
\(n\in\N\), \(x_0=x\), \(x_n=y\), and
\[
 \text{for every }i<n,\qquad x_iRx_{i+1}\ \text{or}\ x_{i+1}Rx_i.
\]
Equivalently each step belongs to \(R\cup R^{-1}\). The sequence is a
function from \(n+1=\{0,\ldots,n\}\) to \(X\), retaining the earlier
finite-sequence convention. A length-zero path has one vertex and no
steps: it connects \(x\) to itself. It is not the length-zero sequence
with no vertex and is not a path between distinct endpoints.
\index{equivalence closure!finite zig-zags}
\index{relation!converse}
\end{hbdefinition}

\begin{hbdefinition}{def:p01-ec-constructions}{Path relation and bounded intersection}
Fix \(X\) and \(R\subseteq X\times X\). Define
\[
 E\defeq\{\langle x,y\rangle\in X\times X:
       \text{a finite zig-zag path along }R\text{ connects }x\text{ to }y\}.
\]
This specifies \(E\) by witnesses, without presupposing that a least
equivalence already exists. It is a set by separation from \(X\times X\).
Finite sequence witnesses can themselves be bounded by
\(\mathcal P(\N\times X)\); they are not a quantification over a
proper collection of paths.

For comparison, let
\[
 \mathcal F\defeq\{F\in\mathcal P(X\times X):
       F\text{ is an equivalence relation on }X\text{ and }R\subseteq F\},
\]
and define the bounded intersection relation
\[
 I\defeq\{p\in X\times X:\text{for every }F\in\mathcal F,
                                      \ p\in F\}.
\]
Here \(F\) ranges over relations, \(p\) over ordered pairs, and
\(\mathcal F\) over a set of such relations. Once nonemptiness of
\(\mathcal F\) is established, \(I=\bigcap\mathcal F\) has the
ordinary nonempty-family intersection interpretation.
\end{hbdefinition}

\begin{hbtheorem}{thm:p01-ec-construction}{proved=annex:p01.ec.construction}{Existence, minimality and the two constructions}
For every set \(X\) and every \(R\subseteq X\times X\), the path
relation \(E\) is reflexive, symmetric and transitive, contains \(R\),
and is contained in every equivalence relation on \(X\) containing \(R\).
Thus it is the unique least such relation.

The family \(\mathcal F\) is nonempty, its intersection \(I\) is an
equivalence relation containing \(R\), and
\[
 E=I=(R\cup R^{-1})^*.
\]
The star in this display is the earlier reflexive-transitive relation
closure, not a set of strings. If \(X=\varnothing\), then
\(R=E=I=\varnothing\). If \(R=\varnothing\) on any \(X\), then
\(E=\{\langle x,x\rangle:x\in X\}\). If \(R\) is already an
equivalence relation on \(X\), then \(E=R\).
\end{hbtheorem}

The symmetry step is essential. Directed paths along \(R\) alone
construct \(R^*\), which need not identify endpoints connected by a
backward step. Nor does merely adding reversed edges ensure transitivity;
the paths may have arbitrarily many, but always finitely many, steps.
No search or termination claim for arbitrary infinite carriers follows
from this existence construction.

\subsection{Classes, representatives and the quotient map}

\begin{hbdefinition}{def:p01-ec-quotient}{Classes and projection on a set}
Let \(E\) now denote any equivalence relation on a set \(A\), as in
the earlier quotient definition. Keep its notation and orientation:
\[
 [a]_E\defeq\{x\in A:xEa\},\qquad
 A/E\defeq\{[a]_E:a\in A\},\qquad q(a)\defeq[a]_E.
\]
The quotient can equivalently be specified as the subset
\[
 \{C\in\mathcal P(A):\text{there exists }a\in A\text{ with }C=[a]_E\}.
\]
Thus each class and the quotient have explicit set bounds. The proposed
projection is \(q:A\to A/E\). The next result establishes its function
and surjection properties and the relation between equality of classes
and the original relation. When applied to the constructed closure,
take \(A=X\) and the \(E\) of \cref{def:p01-ec-constructions}.
\index{quotient set!projection}
\end{hbdefinition}

\begin{hbtheorem}{thm:p01-ec-classes}{proved=annex:p01.ec.classes}{Equality of classes and a genuine quotient}
For \(a,a'\in A\), each class is nonempty, and
\[
 aEa'\quad\Longleftrightarrow\quad[a]_E=[a']_E
 \quad\Longleftrightarrow\quad[a]_E\cap[a']_E\ne\varnothing.
\]
The classes cover \(A\), and different classes are disjoint; consequently
\(A/E\) is a partition into nonempty classes. The rule \(q(a)=[a]_E\)
defines a total surjection, and \(q(a)=q(a')\) exactly when \(aEa'\).
For \(A=\varnothing\), the quotient and the projection's graph are
empty; the cover and surjection assertions still hold.
\end{hbtheorem}

\begin{hbconvention}{def:p01-ec-composition}{Typed function composition}
For a set \(B\) and a given function \(h:A/E\to B\), function
composition means \((h\circ q)(a)=h(q(a))\) for \(a\in A\).
The symbol \(\circ\) here is composition of these typed functions,
not a metavariable for a binary formula connective.
\end{hbconvention}

\begin{hbtheorem}{thm:p01-ec-factorization}{proved=annex:p01.ec.factorization}{Representative independence and the factorization property}
Let \(E\) be an equivalence relation on a set \(A\), let
\(q:A\to A/E\) be the projection, and let \(f:A\to B\) be a
function to any set \(B\). There is a function \(\bar f:A/E\to B\)
satisfying \(f=\bar f\circ q\) if and only if
\[
 \text{for all }a,a'\in A,\qquad aEa'\ \Longrightarrow\ f(a)=f(a').
\]
When it exists, \(\bar f\) is unique and satisfies
\(\bar f([a]_E)=f(a)\). Its graph is the explicitly bounded relation
\[
 \bar f=\{\langle C,b\rangle\in(A/E)\times B:
                \text{there exists }a\in C\text{ with }f(a)=b\}.
\]
This construction does not select a representative in each class.

In particular, if \(E\) is the equivalence closure of
\(R\subseteq X\times X\) and \(f:X\to B\), factorization through
\(X/E\) is equivalent to the simpler generator condition
\[
 \text{for all }x,y\in X,\qquad xRy\ \Longrightarrow\ f(x)=f(y).
\]
Equal values on the generating edges suffice for all finite zig-zags.
\index{quotient set!universal property}
\index{well-definedness!independence of representative}
\end{hbtheorem}

The property is universal for this set-function problem: every
function constant on the identified elements factors through the quotient
in exactly one way. It does not assert that an arbitrary operation or
relation can be applied to representatives without an analogous compatibility
argument. Surjectivity of \(q\) gives uniqueness of the factor; it is not
used to choose an inverse or a section of \(q\). Local use of an existing
member of one specified class in a proof does not choose members of all
classes simultaneously. The full graph proof therefore requires no axiom
of choice.

\subsection{A connected finite construction and a failed quotient formula}

\begin{hbexample}{def:p01-ec-finite-data}{Generators, two candidate blocks and a function}
Let
\[
 X=\{0,1,2,3,4,5\},\qquad
 R=\{\langle0,1\rangle,\langle2,1\rangle,
      \langle2,3\rangle,\langle4,5\rangle\}.
\]
Define the subsets \(C_0=\{0,1,2,3\}\) and \(C_1=\{4,5\}\),
and fix two distinct target elements \(u,v\), with \(B=\{u,v\}\).
Define \(f:X\to B\) by \(f(x)=u\) for \(x\in C_0\) and
\(f(x)=v\) for \(x\in C_1\). These subsets are a disjoint cover
of the displayed finite \(X\), so exactly one clause applies to each
input. Their status as the equivalence classes is established next, rather
than assumed from the appearance of the generators.
\end{hbexample}

\begin{hbtheorem}{thm:p01-ec-finite-example}{proved=annex:p01.ec.finite-example}{The closure, quotient and failure in the finite example}
For the displayed \(X,R\), their equivalence closure is
\[
 E=(C_0\times C_0)\cup(C_1\times C_1),\qquad
 X/E=\{C_0,C_1\}.
\]
The closure contains twenty ordered pairs. Its quotient projection sends
\(0,1,2,3\) to \(C_0\) and \(4,5\) to \(C_1\). The function
\(f\) factors uniquely, with \(\bar f(C_0)=u\) and
\(\bar f(C_1)=v\). In contrast, \(0E3\) but not \(0R^*3\),
so directed reachability would give the wrong quotient identifications.

The identity-valued rule \(g(x)=x\), from \(X\) to \(X\), cannot
descend by \(\bar g([x]_E)=x\): the same class would acquire both
the values zero and one. In this particular finite ordered example one
can instead define \(t:X/E\to X\) by \(t(C)=\min C\).
Here minimum uses the usual order of the natural numbers.
Then \(t(C_0)=0\) and
\(t(C_1)=4\), and \(q\circ t\) is the identity on \(X/E\).
Nevertheless \(t\circ q\ne g\). This is a different, explicitly
specified representative-selection map, not a repair of the failed
factorization of \(g\).
\end{hbtheorem}

The actual zig-zag \(0,1,2,3\) uses a forward step, a backward step,
then a forward step. It explains why \(0\) and \(3\) must be identified
even though the arrows of \(R\) give no directed route from zero to three.
The equality \(f(0)=f(1)=f(2)=f(3)\) makes any representative of that
class safe for \(f\); the unequal values of \(g\) do not. A rule that
chooses a smallest representative is justified here by this finite subset
of \(\N\), not by an assumed canonical representative in every quotient.

\begin{hbmethod}{def:p01-ec-use}{Construct a closure, then descend a function}
Fix the carrier and generators before comparing relations. Permit backward
as well as forward generator steps, and include zero-step paths. Prove
reflexivity, symmetry, transitivity and minimality for the resulting set
relation; the word ``least'' does not supply these proofs.

Form classes only after equivalence is established. To descend a function,
prove that it has equal values on equivalent inputs. For a generated
equivalence, it is enough to check the original generating pairs and then
use the path argument. Construct the quotient function by its graph rather
than assuming a simultaneous choice of representatives. Reject a formula
that gives different outputs on one class, even if a preferred representative
makes a separately defined function possible.
\end{hbmethod}
